\section{set theory}

In the preceding paragraphs, we have seen that predicates can be combined and
quantified. But what are the elementary predicates that we can consider to found
all of mathematics? Set theory provides us with the formalization of the concept
of \emph{membership}: the basic predicate (the only one we will need) is a
predicate of the form $x \in A$, which means "the object $x$ is an element of
the set $A$". Set theory does not explain (nor does it define) what objects and
sets are because these are primitive concepts (just as \emph{points} and
\emph{lines} are for Euclidean geometry). Set theory simply provides us with the
formal rules that can be used to handle predicates of the form $x\in A$. For
example, an \emph{axiom} of set theory is the following.

\begin{axiom}[empty set]
\[
  \exists A \colon \lnot \exists x\colon x \in A
\]
\end{axiom}

The axiom means: "there exists a set $A$ that contains no element $x", that is,
we are saying that the \emph{empty set}%
\mymargin{empty set}%
\index{empty set} exists, which is usually denoted by the symbol $\emptyset$. We will thus write $\exists\emptyset\colon \forall x\colon \lnot(x\in \emptyset)$. We can then introduce a demonstrative context in which $\emptyset$ is a constant, and the axiom becomes more simply $\forall x\colon \lnot (x\in \emptyset)$.

Other appropriate axioms of set theory guarantee the existence of the union
$A\cup B$, intersection $A\cap B$, and difference $A\setminus B$ of any two sets
$A$ and $B$. These operations between sets can be traced back to the membership
relation and can thus be defined as follows.

\begin{axiom}[operations between sets]
If $A$ and $B$ are sets, then the sets $A\cup B$ (union), $A\cap B$
(intersection), and $A\setminus B$ (difference) exist.
\mymargin{union intersection difference}%
\index{union}%
\index{intersection}%
\index{difference}%
Formally:
\begin{align*}
    \forall A,B\colon \exists A\cup B \colon \forall x\colon x\in A\cup B &\iff (x\in A) \lor (x\in B),\\
    \forall A,B\colon \exists A\cap B \colon \forall x\colon x\in A\cap B &\iff (x\in A) \land (x\in B),\\
    \forall A,B\colon \exists A\setminus B \colon \forall x\colon x\in A\setminus B &\iff (x\in A) \land \lnot (x \in B).
\end{align*}
\end{axiom}

The negation of membership $\lnot (x \in A)$ is usually abbreviated as $x \not
\in A$.

Using the simple membership relation, we can define the relations of inclusion
and equality between sets: $A \subset B$ is read as "$A$ is a subset of $B$", $A
\supset B$ is read as "$A$ is a superset of $B$", and $A=B$ is read as "$A$ is
equal to $B$". These relations are defined by the following properties:
\begin{align*}
  A \subset B &\iff \forall x\colon (x\in A \implies x\in B)\\
  A \supset B &\iff \forall x\colon (x\in A \implied x\in B)\\
  A = B &\iff \forall x\colon (x\in A \iff x \in B).
\end{align*}
In words, we will say that $A$ is a subset of $B$, $A \subset B$, if every
element of $A$ is also an element of $B$. We will say that $A$ and $B$ are
equal, $A=B$, if $A$ and $B$ have the same elements. Obviously, the relation
$\supset$ is nothing but the inverse relation of $\subset$, that is, $A\supset
B$ if and only if $B\subset A$. We will also write $A \neq B$ to indicate the
opposite relation of equality, that is: $\lnot(A=B)$. Note that $A=B$ is
equivalent to $(A\subset B) \land (B\subset A)$.

It is trivial to observe that $\forall A\colon A\subset A$ since $x\in A
\implies x\in A$. Thus, the axiom of equality, $A=A$, can actually be proven.

The relation $A \subset B$ means that every element of $A$ is also an element of
$B$, and therefore there are no elements of $A$ that are not in $B$, that is,
$A$ is not larger than $B$. As a mnemonic device, note that the symbol $\subset$
is oriented so that the smaller set is on the narrower side of the relation,
just as in the inequality $3 \le 5$, the smaller number is on the narrower side
of the symbol $\le$. To prove that the relation $A \subset B$ is true, one must
take any element of $A$ and prove that this element is also an element of $B$.
To prove that the equality $A=B$ holds, one typically proceeds by separately
proving the two inclusions $A\subset B$ and $B\subset A$.

The ability to construct sets of sets is very useful. For this reason, set
theory usually makes no distinction between objects and sets of objects. In the
primitive relation $x\in A$, $x$ can also be a set. We can then imagine that
every object in our universe is a set. In this way, the equality relation $A=B$
that we have defined above is well-defined for every pair of objects (or sets,
which is the same) $A$ and $B$.

Since the elements of a set $\mathcal A$ are themselves sets, it is possible to
consider the union $\bigcup \mathcal A$ and, if\mynote{%
The intersection of an empty family of sets would yield the universal set, which
we will see cannot be defined.
} 
$\mathcal A \neq \emptyset$, the intersection $\bigcap \mathcal A$ of all the
elements of $\mathcal A$. This extends the concept of union and intersection
even to possibly infinite families.

\begin{axiom}[arbitrary union and intersection]
If $\mathcal A$ is any set, the set $\bigcup \mathcal A$ exists, and if
$\mathcal A \neq \emptyset$, the set $\bigcap \mathcal A$ exists with the
following properties:
\begin{align*}
  x \in \bigcup \mathcal A & \iff \exists A \in \mathcal A \colon x\in A, \\
  x \in \bigcap \mathcal A & \iff \forall A \in \mathcal A \colon x\in A.
\end{align*}
\end{axiom}
An alternative notation is the following:
\[
 \bigcup \mathcal A = \bigcup_{A\in \mathcal A} A, \qquad 
 \bigcap \mathcal A = \bigcap_{A\in \mathcal A} A.  
\]
This notation highlights the fact that the set $\bigcup \mathcal A$ is the union
of all the elements $A$ of the set $\mathcal A$, while $\bigcap \mathcal A$ is
the intersection of all the elements of $\mathcal A$. The set $\mathcal A$ is
often called a \emph{family} of sets because its elements are treated as sets
rather than objects.

Another axiom of set theory guarantees that for every $x$, there exists a set
whose only element is $x$.
\begin{axiom}[singleton]
  For every $x$, there exists the set $\ENCLOSE{x}$, called a \emph{singleton}%
\mymargin{singleton}%
\index{singleton}
  such that:
  \[
    y \in \ENCLOSE{x} \iff y=x.
  \]
\end{axiom}
By taking the union of singletons, we can define (by enumeration) sets that
contain a finite number of objects:
\begin{align*}
  \ENCLOSE{a,b} &= \ENCLOSE{a} \cup \ENCLOSE{b} \\
  \ENCLOSE{a, b, c} &= \ENCLOSE{a,b} \cup \ENCLOSE{c}\\
  \ENCLOSE{a, b, c, d} &= \ENCLOSE{a,b,c} \cup \ENCLOSE{d}\\
  &\quad\vdots
\end{align*}

However, be careful: the set $\ENCLOSE{a,b}$ contains two elements only if
$a\neq b$, in fact, if $a=b$, it can be easily verified that
\begin{equation}\label{eq:4775523}
\ENCLOSE{a,a} = \ENCLOSE{a}.
\end{equation}

\begin{exercise}
  Verify \eqref{eq:4775523} using the formal definitions given earlier.
\end{exercise}
\begin{proof}[Solution.]
Using the definitions of union, singleton, and logical disjunction, we have the
equivalence of the following predicates:
\begin{gather*}
  x \in \ENCLOSE{a,a}  \\
  x \in \ENCLOSE{a} \cup \ENCLOSE{a}\\
  (x \in \ENCLOSE{a}) \lor (x \in \ENCLOSE{a})\\
  (x = a) \lor (x = a) \\
  x = a \\
  x \in \ENCLOSE{a}
\end{gather*}
and thus $\ENCLOSE{a,a}=\ENCLOSE{a}$ by the definition of equality between sets.
\end{proof}

If $P(x)$ is a predicate in a single variable $x$, we would like to be able to
define the set of all objects $x$ that make the predicate $P$ true.
Surprisingly, if we were to add this axiom (called the naive specification
axiom) to set theory, we would encounter a paradox%
\index{naive set theory}%
\index{sets!naive theory}%
\index{Cantor, Georg}%
\index{Frege}%
\index{Russell}%
\mynote{Georg Cantor (1845--1918), Bertrand Russell (1872--1970) see historical notes on page~\pageref{nota:Cantor}}.

\begin{theorem}[Russell's paradox]
\label{th:Russell}%
There does not exist a set $R$ such that for every $x$, we have
\[
   x\in R \iff x\not \in x
\]
(such a set could be defined in this way through a \emph{naive} specification
axiom: $R=\ENCLOSE{x\colon x\not \in x}$).
\end{theorem}
%
\begin{proof}
  If it existed, we would have:
  \[
    R \in R 
    \iff R\not \in R
  \]
    and this contradicts the principle of non-contradiction (a predicate and its
  negation cannot both be true).
\end{proof}

To avoid inconsistency, it is necessary to limit the \emph{specification axiom}
to the construction of subsets of already constructed sets.

\begin{axiom}[specification]
    If $P$ is any predicate with a free variable $x$ and if $B$ is a set, then
  there exists the set $\ENCLOSE {x\in B\colon P(x)}$ formed by all the elements
  of $B$ that satisfy the predicate $P$:
\[
  a \in \ENCLOSE{x\in B\colon P(x)} \iff (a \in B \land P(a)).
\]
\end{axiom}

If we now try to repeat Russell's reasoning given a set $A$, we can construct
the set $R=\ENCLOSE{x\in A\colon x\not \in x}$ and discover that it cannot be
$R\in R$ (otherwise it should be $R\not \in R$, absurd). But it could very well
be $R\not \in R$ if also $R\not \in A$. Thus, it follows that given any set $A$,
there is a set that is not an element of $A$: we must therefore resign ourselves
to the fact that there is no \emph{universal} set containing every other set,
but at least we do not obtain a contradiction%
\mynote{%
On the other hand, we cannot exclude that a contradiction exists. 
In fact, Goedel has shown that it is not possible to prove that there are no
contradictions in the axiomatic system we are considering.
}.

To complete set theory, we also introduce the concept of the \emph{power set}%
\mymargin{power set}%
\index{set!power set}
\index{$\mathcal P(\cdot)$}% 
that is, the set of all subsets of a given set.
\begin{axiom}[power set]
\label{def:insieme_parti}%
If $A$ is a set, the set $\P(A)$ of the parts of $A$ exists, which is the set of
all subsets of $A$:
\begin{equation}\label{eq:insieme_delle_parti}
  x \in \mathcal P(A) \iff x \subset A.
\end{equation}
\end{axiom}

\subsection{Von Neumann's finite ordinals}

With the only axioms we have introduced so far, we can construct an arbitrarily
large number of finite sets through a construction due to Von Neumann. We can
define
\begin{equation}\label{eq:vonNeumann}
  \begin{aligned}
  0&=\emptyset, &
  1&=0\cup\ENCLOSE{0}, &
  2&=1\cup\ENCLOSE{1}, &
  3&=2\cup\ENCLOSE{2}, &
  4&=3\cup\ENCLOSE{3},\\
  5&=4\cup\ENCLOSE{4}, &
  6&=5\cup\ENCLOSE{5}, &
  7&=6\cup\ENCLOSE{6}, &
  8&=7\cup\ENCLOSE{7}, &
  9&=8\cup\ENCLOSE{8}.
  \end{aligned}
\end{equation}
By making these definitions explicit, we observe that:
$0=\ENCLOSE{}$ is a set with zero elements, 
$1=\ENCLOSE{0}=\ENCLOSE{\ENCLOSE{}}$ is a set with one element (zero),
$2=\ENCLOSE{0,1}=\ENCLOSE{\ENCLOSE{},\ENCLOSE{\ENCLOSE{}}}$ is a set with two
elements (the two preceding ones),
$3=\ENCLOSE{0,1,2}=\ENCLOSE{\ENCLOSE{},\ENCLOSE{\ENCLOSE{}},\ENCLOSE{\ENCLOSE{},\ENCLOSE{\ENCLOSE{}}}}$
is a set with three elements (the three preceding ones),
etc.
Proceeding at will, it is possible to define an arbitrarily large number of
finite sets, each different from the other.

These particular sets certainly exist, and we can therefore use them in the
following to make examples of finite sets. It is also possible (but not
necessary) to use these sets to define natural numbers (see the Axiom of
Infinity~\ref{axiom:infinito}).

\begin{exercise}
Verify that $4\in 5$, 
that $4\subset 5$, 
and that $4\neq 5$.
\end{exercise}

\subsection{relations}

Thanks to the previous axioms, it is possible to define the Cartesian product of
two sets: this will be a very important concept in the following. First, we must
define the concept of a \emph{pair}%
\mymargin{pair}%
\index{pair}: given
$a,b$, we define the pair $(a,b)$ with the first element $a$ and the second
element $b$ as an object that has this property:%
\mynote{%
A formal way to define the pair using sets is due to Kuratowski:
\[
   (a,b) = \ENCLOSE{\ENCLOSE{a},\ENCLOSE{a,b}}.
\]
It is not difficult to verify that this definition satisfies the 
property~\eqref{eq:coppia}.
Thus, we can define the product set
\begin{align*}
  A\times B = \big\{&C\in \mathcal P(\mathcal P(A\cup B))\colon \\ 
  & \exists a\colon \exists b\colon a\in A, b\in B, \\
  & C=\ENCLOSE{\ENCLOSE{a},\ENCLOSE{a,b}}\big\}
\end{align*}
}%
\begin{equation}\label{eq:coppia}
  (a, b) = (a', b') \iff (a=a') \land (b=b').
\end{equation}
We are thus requiring that a pair be identified by the two elements that compose
it, in which, however, the order in which they are listed is important (unlike
the set $\ENCLOSE{a,b}$, in which the order of elements is irrelevant). In the
following, we will also use the notation $a \mapsto b$ to indicate the pair
$(a,b)$ and will call it an \emph{arrow}%
\mymargin{arrow}%
\index{arrow} from $a$ to $b$.

The \emph{Cartesian product} $A\times B$ of two sets $A$ and $B$ is the set of
all pairs whose first element is in $A$ and the second element is in $B$:
\[
  (a, b) \in A \times B \iff (a\in A \land b\in B).
\]

If we represent the elements of $A$ as points on a horizontal line (the
$x$-axis) and the elements of $B$ as points on a vertical line (the $y$-axis),
the elements of $A\times B$ can be represented as points in the plane that
project onto $A$ on the $x$-axis and onto $B$ on the $y$-axis (see
figure~\ref{fig:funzione}).
  
A \emph{relation}%
\mymargin{relation}%
\index{relation} $R$ between the elements of a set $A$ and the elements of a set $B$ is nothing but a subset of $A\times B$, that is, $R\in \mathcal P(A\times B)$. For a relation $R\subset A\times B$, the infix notation $aRb$ is used to indicate $(a,b)\in R$. Alternatively, with the arrow notation, we can write $a \stackrel R \mapsto b$ to indicate that the arrow $a\mapsto b$ is an element of $R$. For example, if we take $A=B=\ENCLOSE{1,2,3}$ and consider the relation $R=\ENCLOSE{(1,2),(1,3),(2,3)}$, the predicate $aRb$ represents the usual order relation $a<b$ on the three numbers considered.

\subsubsection{equivalence relations}

\begin{definition}[equivalence relation]
\label{def:equivalenza}%
\label{def:insieme_quoziente}%
Let $R\subset A\times A$ be a relation on the set $A$. We will say that $R$ is
an \emph{equivalence relation}%
\mymargin{equivalence relation}%
\index{relation!equivalence} if the following properties (typical of equality) hold for every $x,y,z\in A$:
\begin{enumerate}
  \item reflexive: $x R x$;
  \item symmetric: if $x R y$, then $y R x$;
  \item transitive: if $x R y$ and $yRz$, then $x R z$.
\end{enumerate}
If $R$ is an equivalence relation, we will say that $x$ is equivalent to $y$
(via $R$) when $xRy$. Given $x \in A$, the set of all elements $y\in A$ that are
equivalent to $x$ is called the \emph{equivalence class}%
\mymargin{equivalence class}%
\index{class!equivalence} of $x$. It is defined in this way:
\[
  [x]_R = \ENCLOSE{y\in A \colon xRy}.  
\]
If $R$ is an equivalence relation on $A$, we define the \emph{quotient set}%
\mymargin{quotient set}%
\index{set!quotient}%
$A/R$
as the set of all equivalence classes:
\[
 A/R 
 = \ENCLOSE{[x]_R\colon x\in A} 
 = \ENCLOSE{B\in \mathcal P(A)\colon \exists x\in A\colon B=[x]_R}.  
\]
\end{definition}

If $R$ is an equivalence relation on $A$, the quotient set $A/R$ represents the
set of objects in $A$ where equivalent objects are identified with each other.
We will often say that $A/R$ represents $A$ modulo the equivalence $R$.

For example, if $A=\ENCLOSE{1,2,3,4,5}$ and we take $D=\ENCLOSE{1,3,5}$ (the odd
numbers) and $P=\ENCLOSE{2,4}=A\setminus D$ (the even numbers), we can define a
relation $R$ on $A$ by the property:
\[
 x R  y \iff (x\in P \land y\in P) \lor (x\in D \land y\in D).
\]
The relation $R$ represents the property of the elements of $A$ of having the
same "parity". We will have $[1]_R = [3]_R=[5]_R= D$ and $[2]_R=[4]_R=P$. Thus
\[
   A/R = \ENCLOSE{P,D} = \ENCLOSE{\ENCLOSE{2,4},\ENCLOSE{1,3,5}}
\]
and, indeed, this set has two elements $P$ and $D$, which is what is obtained by
identifying the numbers of $A$ based on the property of being even or odd.

The quotient set $A/R$ is a \emph{partition}%
\mymargin{partition}%
\index{partition} of $A$ because it is formed by disjoint sets whose union is all of $A$. In fact, giving a partition of $A$ is equivalent to giving an equivalence relation.

\subsubsection{order relations}

\begin{definition}[order relation]
  \label{def:ordine}%
  A relation
  $\le$ on a set $X$ is said to be an
  \emph{order relation}%
\mymargin{order relation}%
\index{relation!order}
  if it satisfies the following properties (for every $x,y,z\in X$):
  \begin{enumerate}
    \item[1.] reflexive: $x\le x$;
    \item[2.] antisymmetric: $x\le y \land y\le x \implies x=y$;
    \item[3.] transitive: $x\le y \land y\le z \implies x\le z$.
  \end{enumerate}
  Furthermore, the order relation $\le$
  is said to be a \emph{total order}%
\mymargin{total order}%
  (or \emph{linear}) 
  \index{ordering!total}%
  \index{ordering!linear}%
  \index{total!order}%
  \index{linear!order}%
  and $X$ is said to be \emph{totally ordered} if
  \index{totally ordered}%
  \begin{enumerate}
    \item[4.] trichotomy: $x\le y \lor y\le x$.
  \end{enumerate}
\end{definition}

\begin{definition}
  \label{def:massimo}%
  \label{def:minimo}%
  \label{def:maggiorante}%
  \label{def:minorante}%
  \label{def:sup}%
  \label{def:inf}%
    Let $\le$ be an order relation on a set $X$, let $A\subset X$, and let $x\in
  X$.

    We will say that $x$ is an \emph{upper bound} of $A$ if for every $a\in A$,
  we have $a\le x$.
  \mymargin{upper bound, lower bound}%
    We will say that $x$ is a \emph{lower bound} of $A$ if for every $a\in A$,
  we have $x\le a$.
  
  We will say that $x$ is the \emph{maximum} of $A$ 
  and write $x=\max A$,
  if $x$ is an upper bound of $A$ and $x\in A$.
  \mymargin{maximum, minimum}
  We will say that $x$ is the \emph{minimum} of $A$ 
  and write $x=\min A$,
  if $x$ is a lower bound of $A$ and $x\in A$.
  
  We will say that $x$ is the \emph{supremum} of $A$ 
  and write $x=\sup A$,
  if $x$ is an upper bound of $A$
  and there are no upper bounds of $A$ smaller than $x$
  ($x$ is the minimum of the upper bounds).
  \mymargin{supremum and infimum}%
  We will say that $x$ is the \emph{infimum} of $A$ 
  and write $x=\inf A$,
  if $x$ is a lower bound of $A$
  and there are no lower bounds of $A$ larger than $x$
  ($x$ is the maximum of the lower bounds).

  If the set $A$ has at least one upper bound $x\in X$, we will say that $A$ 
  \mymargin{bounded}%
  \index{bounded}%
  \index{set!bounded}%
  is \emph{bounded above}. 
    If $A$ has at least one lower bound, we will say that $A$ is \emph{bounded
  below}.
    If $A$ is both bounded above and bounded below, we will say that $A$ is
  \emph{bounded}.
\end{definition}

Note that the maximum and minimum of a set $A$, if they exist, are unique. In
fact, if $x$ and $y$ are both maxima of $A$, we must have $x\le y$ since $y$ is
a maximum, but also $y\le x$ since $x$ is also a maximum. Thus, $x=y$ by the
antisymmetric property. The same holds for the minimum. Since the supremum is
the minimum of the upper bounds, and the infimum is the maximum of the lower
bounds, the supremum and infimum, if they exist, are also unique. If the maximum
(or minimum) exists, then it coincides with the supremum (or infimum).
Conversely, we will see examples where the supremum exists but is not an element
of the set (thus it is not a maximum, and the maximum does not exist).

In general, if $\le$ is an order relation on $X$, the inverse relation $\ge$ is
defined by setting $x\ge y$ when $y\le x$. The reflexive property identifies
\emph{wide} orderings. The corresponding strict ordering $<$ is defined by
setting $x < y$ when $x\le y \land x\neq y$, and the inverse relation $>$ is
defined by setting $x>y$ when $y<x$.

\begin{example} 
  Let $X=\ENCLOSE{1,2,3,4}$. 
    It can be verified that there is a unique total order relation on $X$ for
  which $1 \le 2$, $2 \le 3$, and $3 \le 4$ (we will write more briefly $1 \le 2
  \le 3 \le 4$).
  This relation is defined by the set of pairs:
  \[
    R_\le = \ENCLOSE{
          (1,1),(1,2),(1,3),(1,4),
          (2,2),(2,3),(2,4),
          (3,3),(3,4),
          (4,4)}
  \]
  in the sense that $x\le y$ if $(x,y)\in R_\le$.
  The inverse relation $\ge$ corresponds to:
  \[
    R_\ge = \ENCLOSE{
          (1,1),(2,1),(3,1),(4,1),
          (2,2),(3,2),(4,2),
          (3,3),(4,3),
          (4,4)}
  \]
  and the strict relations $<$ and $>$ correspond to:
  \begin{align*}
    R_< &= \ENCLOSE{
          (1,2),(1,3),(1,4),
          (2,3),(2,4),
          (3,4)}, \\
    R_> &= \ENCLOSE{
          (2,1),(3,1),(4,1),
          (3,2),(4,2),
          (4,3)}.
  \end{align*}

  Taking $A=\ENCLOSE{1,3}\subset X$, we have $\max A = \sup A = 3$ and 
  $\min A = \inf A = 1$.
  The upper bounds of $A$ are $3$ and $4$. 
  There is a unique lower bound, which is $1$.
\end{example}

\begin{example}
If we take the set $A=\ENCLOSE{1,2,3}$ and consider the relation on $X=\mathcal
P(A)$ defined by $x\le y$ if $x\subset y$, we obtain an order relation that is
not total.
In fact, if $x=\ENCLOSE{1,2}$ and $y=\ENCLOSE{2,3}$, it is neither true that
$x\subset y$ nor that $y\subset x$.
We have $\max \P(A) = \sup \P(A) = A$ and $\min \P(A) = \inf \P(A) = \emptyset$.
\end{example}

\begin{example}
The set $\NN$ of natural numbers (introduced in chapter~\ref{sec:naturali}) has
no upper bounds, so it has neither a maximum nor a supremum.
In fact, given any natural number, there is always another larger one.

If we add an element $\infty$ to the set of natural numbers, which we define to
be larger than any natural number, we obtain a totally ordered set
$X=\NN\cup\ENCLOSE{\infty}\supset \NN$ such that $\max X = \infty$ and $\sup \NN
= \infty$ in $X$. But $\infty$ is not an element of $\NN$, so it is not a
maximum of $\NN$ in $X$.
\end{example}

\begin{definition}[dense ordering]
  \label{def:ordinamento_denso}%
  An ordering is said to be
  \emph{dense}%
\mymargin{dense ordering}%
\index{dense} or \emph{divisible} 
  if between two distinct points there is always an intermediate point:
  \index{dense!subset}%
  \index{dense!ordering}%
  \index{subset!dense}%
  \index{ordering!dense}%
  \mynote{More generally, if $X$ is a totally ordered set and $A\subset X$, $A$ is said to be dense in $X$ if given any $x,y\in X$ with $x<y$, there exists $c\in A$ such that $x<c<y$.}
  \[
   x < y \implies \exists c \colon x < c < y.
  \]
\end{definition}

\begin{example}
An example of a dense ordering is that given on the set of points of a geometric
line. Given two distinct points $P$ and $Q$ on the line, there is always an
intermediate point between $P$ and $Q$ (for example, the midpoint).

Intuitively, we observe that if we have a dense ordering on a set and the set
contains at least two points, then the set contains infinitely many points since
there is an intermediate point, and then another intermediate point between the
intermediate point and one of the two points, and so on...
\end{example}

\begin{definition}[continuous ordering]
  \label{def:ordinamento_continuo}%
  A dense ordering $\le$ on a set $X$ is said to be
  \emph{continuous}%
\mymargin{continuous ordering}%
\index{continuous}
  \index{continuity!ordering}%
  \index{order!continuous}%
  (or \emph{Dedekind-complete})
  \index{Dedekind-complete}%
  \index{complete!Dedekind}%
  if given any non-empty subsets $A,B$ of $X$
    such that for every $a\in A$ and every $b\in B$, we have $a\le b$
    (concisely, we will write $A\le B$ and say that $A$ and $B$ are 
    \emph{separated})
    then there exists $c\in X$ such that for every $a\in A$ and every $b\in B$, 
    we have $a\le c \le b$ (concisely, we will write $A\le c \le B$
    and say that $c$ is a separation element between $A$ and $B$).
\end{definition}

In words, the previous definition is expressed by saying that a densely ordered
set is \emph{continuous} if every pair of separated subsets admits a separation
element. In many texts, this condition is called \emph{completeness} or, more
precisely, \emph{Dedekind completeness}.
\index{completeness!Dedekind}%
\index{Dedekind!complete}%

\begin{example}
A set can be dense but not continuous. 
For example, if we take two distinct points $P$, $Q$ on an oriented line, such
that $P<Q$, then take the midpoint $M$ between $P$ and $Q$, then the midpoint
between all the previous points, and so on, intuitively, we obtain a set of
points $X$ that turns out to be dense. But if we take the set $A$ of points of
$X$ that are less than twice as far from $P$ as they are from $Q$, and the set
$B=X\setminus A$, it is clear that the two sets are non-empty and separated, but
there is no separation point between $A$ and $B$ because such a point should be
at a distance of $\frac 1 3 PQ$ from $P$, while all the points of $X$ have a
distance from $P$ that is expressed as a fraction of $PQ$ with a power of $2$ as
the denominator.
\end{example}

\mynote{
We will avoid using the definition of continuity on orderings that are not
dense; otherwise, all finite sets would have a continuous ordering, and even
sets like $\NN$ and $\ZZ$ would result in being continuous, which goes against
the intuition of continuity we have in mind.
}

\begin{theorem}[existence of the $\sup$]%
  \label{th:sup}%
  \mymark{**}%
  Let $X$ be a set with a continuous ordering.
  If $A\subset X$ is a non-empty set
  and bounded above, then $\sup A$ (the supremum) exists.
  If $A\subset X$ is a non-empty set and bounded below, 
  then $\inf A$ (the infimum) exists.

  If $X$ has a minimum, then $\sup\emptyset = \min X$
  and if $X$ has a maximum, then $\inf\emptyset = \max X$.
  \end{theorem}
  %
  \begin{proof}
  \mymark{*}
  Consider the set of upper bounds of $A$:
  \[
  B = \ENCLOSE{ b\in \RR \colon b \ge A}.
  \]
  Given $a\in A$ and $b\in B$, since $b$ is an upper bound of $A$,
  we must have $a\le b$. Thus, $A$ and $B$ are separated.
  Moreover, $A$ is non-empty by hypothesis, and $B$ is non-empty because 
  $A$ is bounded above, and thus there exists at least one upper bound of $A$.
  Therefore, since the ordering is continuous by hypothesis,
    there must exist $c\in X$ that is an upper bound of $A$ and a lower bound of
  $B$.
  Since $c$ is an upper bound of $A$, we have $c\in B$.
  But $c$ is also a lower bound of $B$, so $c=\min B$.
  Being $c$ the minimum of the upper bounds of $A$, it is, by definition,
  $c= \sup A$.

  An analogous proof holds for the infimum.

  As for the empty set, note that every $x\in X$ 
  is both an upper bound and a lower bound of the empty set.
  Thus, the set of upper bounds and the set of lower bounds coincide with $X$.
    The minimum of the upper bounds is thus $\min X$, and the maximum of the
  lower bounds is $\max X$.
\end{proof}

We observe that the existence of the $\sup$
of non-empty and bounded above sets
is indeed a condition equivalent to the axiom of continuity.
In fact, if $A\le B$, certainly $\sup A$, if it exists, is an element 
of separation between $A$ and $B$.

The geometric line is an example of a set that 
we would like to have the property of continuity.
This is because when a line is drawn that crosses a geometric figure, 
we would like to have the certainty that, 
if the line has a point inside the figure and a point outside, 
then it also has a point of intersection (for us, this will be the theorem 
of zeros~\ref{th:zeri}).
These points of intersection can be  
seen as points of separation between the points
inside and the points outside the figure:
for this reason, we are interested in the properties of continuity.

The issue is delicate and not trivial. 
Instead of aiming to formalize the axioms of Euclidean geometry,
we will have a more algebraic approach in the sense that we will define 
the geometric line using the set $\RR$ of real numbers.
Much of this chapter will be essentially dedicated to this purpose.
But see the discussion on Euclid and Hilbert in the historical notes 
at the end of the chapter.